\documentclass[10pt,a4paper]{amsart}
\usepackage{amsmath,amssymb,amsthm,graphicx}
% For commutative diagrams we use Xy-pic
\usepackage[all]{xy}

% We use 2cell for 2-commutative diagrams.
\xyoption{2cell}
% Theorem environments.
%
\theoremstyle{plain}
\newtheorem{theorem}[subsection]{Theorem}
\newtheorem{proposition}[subsection]{Proposition}
\newtheorem{lemma}[subsection]{Lemma}

\theoremstyle{definition}
\newtheorem{definition}[subsection]{Definition}
\newtheorem{example}[subsection]{Example}
\newtheorem{exercise}[subsection]{Exercise}
\newtheorem{situation}[subsection]{Situation}

\theoremstyle{remark}
\newtheorem{remark}[subsection]{Remark}
\newtheorem{remarks}[subsection]{Remarks}

\numberwithin{equation}{subsection}

% Macros
%
\def\lim{\mathop{\mathrm{lim}}\nolimits}
\def\colim{\mathop{\mathrm{colim}}\nolimits}
\def\Spec{\mathop{\mathrm{Spec}}}
\def\Hom{\mathop{\mathrm{Hom}}\nolimits}
\def\Ext{\mathop{\mathrm{Ext}}\nolimits}
\def\SheafHom{\mathop{\mathcal{H}\!\mathit{om}}\nolimits}
\def\SheafExt{\mathop{\mathcal{E}\!\mathit{xt}}\nolimits}
\def\Sch{\mathit{Sch}}
\def\Mor{\mathop{\mathrm{Mor}}\nolimits}
\def\Ob{\mathop{\mathrm{Ob}}\nolimits}
\def\Sh{\mathop{\mathit{Sh}}\nolimits}
\def\NL{\mathop{N\!L}\nolimits}
\def\CH{\mathop{\mathrm{CH}}\nolimits}
\def\proetale{{pro\text{-}\acute{e}tale}}
\def\etale{{\acute{e}tale}}
\def\QCoh{\mathit{QCoh}}
\def\Ker{\mathop{\mathrm{Ker}}}
\def\Im{\mathop{\mathrm{Im}}}
\def\Coker{\mathop{\mathrm{Coker}}}
\def\Coim{\mathop{\mathrm{Coim}}}

% Boxtimes
%
\DeclareMathSymbol{\boxtimes}{\mathbin}{AMSa}{"02}

%
% Macros for moduli stacks/spaces
%
\def\QCohstack{\mathcal{QC}\!\mathit{oh}}
\def\Cohstack{\mathcal{C}\!\mathit{oh}}
\def\Spacesstack{\mathcal{S}\!\mathit{paces}}
\def\Quotfunctor{\mathrm{Quot}}
\def\Hilbfunctor{\mathrm{Hilb}}
\def\Curvesstack{\mathcal{C}\!\mathit{urves}}
\def\Polarizedstack{\mathcal{P}\!\mathit{olarized}}
\def\Complexesstack{\mathcal{C}\!\mathit{omplexes}}
% \Pic is the operator that assigns to X its picard group, usage \Pic(X)
% \Picardstack_{X/B} denotes the Picard stack of X over B
% \Picardfunctor_{X/B} denotes the Picard functor of X over B
\def\Pic{\mathop{\mathrm{Pic}}\nolimits}
\def\Picardstack{\mathcal{P}\!\mathit{ic}}
\def\Picardfunctor{\mathrm{Pic}}
\def\Deformationcategory{\mathcal{D}\!\mathit{ef}}
\usepackage{cleveref}
\crefname{originalreference}{original reference}{original references}
\Crefname{originalreference}{Original reference}{Original references}
\makeatletter
\newlabel{item-surjective}{{06AP}{source}{}{sourcepointer-1}{}}
\newlabel{item-surjective@cref}{{[originalreference][0][]06AP}{[1][0][]source}{}{}{}}
\makeatother
\crefname{originalreference}{original reference}{original references}
\Crefname{originalreference}{Original reference}{Original references}
\makeatletter
\newlabel{item-vanishing}{{0A53}{source}{}{sourcepointer-2}{}}
\newlabel{item-vanishing@cref}{{[originalreference][0][]0A53}{[1][0][]source}{}{}{}}
\makeatother
\crefname{originalreference}{original reference}{original references}
\Crefname{originalreference}{Original reference}{Original references}
\makeatletter
\newlabel{lemma-even-easier}{{03SG}{source}{}{sourcepointer-3}{}}
\newlabel{lemma-even-easier@cref}{{[originalreference][0][]03SG}{[1][0][]source}{}{}{}}
\makeatother
\crefname{originalreference}{original reference}{original references}
\Crefname{originalreference}{Original reference}{Original references}
\makeatletter
\newlabel{lemma-finite-pushforward-statements}{{0A5D}{source}{}{sourcepointer-4}{}}
\newlabel{lemma-finite-pushforward-statements@cref}{{[originalreference][0][]0A5D}{[1][0][]source}{}{}{}}
\makeatother
\crefname{originalreference}{original reference}{original references}
\Crefname{originalreference}{Original reference}{Original references}
\makeatletter
\newlabel{lemma-restrict-to-open}{{02RX}{source}{}{sourcepointer-5}{}}
\newlabel{lemma-restrict-to-open@cref}{{[originalreference][0][]02RX}{[1][0][]source}{}{}{}}
\makeatother
\crefname{originalreference}{original reference}{original references}
\Crefname{originalreference}{Original reference}{Original references}
\makeatletter
\newlabel{situation-what-to-prove}{{0A52}{source}{}{sourcepointer-6}{}}
\newlabel{situation-what-to-prove@cref}{{[originalreference][0][]0A52}{[1][0][]source}{}{}{}}
\makeatother
\crefname{originallemma}{original lemma}{original lemmas}
\Crefname{originallemma}{Original lemma}{Original lemmas}
\makeatletter
\newlabel{varieties-lemma-normalization-locally-algebraic}{{0BXR}{source}{}{sourcepointer-7}{}}
\newlabel{varieties-lemma-normalization-locally-algebraic@cref}{{[originallemma][0][]0BXR}{[1][0][]source}{}{}{}}
\makeatother
\crefname{originallemma}{original lemma}{original lemmas}
\Crefname{originallemma}{Original lemma}{Original lemmas}
\makeatletter
\newlabel{varieties-lemma-regular-point-on-curve}{{0B8Y}{source}{}{sourcepointer-8}{}}
\newlabel{varieties-lemma-regular-point-on-curve@cref}{{[originallemma][0][]0B8Y}{[1][0][]source}{}{}{}}
\makeatother
\title{The vanishing and surjectivity statements hold for the extension by zero of a constant sheaf along an open of a reduced curve}
\author{The Stacks Project Authors}
\date{stacks-project:etale-cohomology.tex:local-snapshot-20261011; GNU FDL 1.2 or later}
\begin{document}
\maketitle
\noindent\textit{Editorial note.} All delivered mathematical text is the author's own, copied from the pinned Stacks Project chapter file etale-cohomology.tex; no statement or proof step was rewritten, reordered or removed.\par\medskip
\noindent\textit{Editorial note.} The counted proof is the author's complete argument: the proof is written out in full: the normalization and the reduction to the smooth case, the pushforward along the finite morphism, the agreement on the dense open, and the two applications of the restriction-to-open lemma.\par\medskip
\noindent\textit{Editorial note.} The label lemma-somewhat-easier names the counted statement, and the proof environment immediately following it in the source is delivered verbatim as its proof.\par\medskip
\noindent\textit{Editorial note.} The counted statement is carried by the author's own lemma environment in the chapter source.\par\medskip
\noindent\textit{Editorial note.} The statement refers to the standing situation and to the enumerated statements by their labels; the eight cited inputs are kept as external pointers with their canonical Stacks tags.\par\medskip
\noindent\textit{Allowed standard prerequisites (editorial).}\begin{itemize}
\item The case of a smooth curve. Same chapter, the companion lemma proving the same statements under the additional assumption that X is smooth. Supplies the result for the normalization of the curve.
\item Finiteness and smoothness of the normalization of a curve. Varieties, the lemmas stating that the normalization morphism of a curve is finite and that the normalization of a reduced curve is smooth. Supplies the finite morphism nu : X\textasciicircum{}nu -> X and the smoothness of X\textasciicircum{}nu.
\item Pushforward of the statements along a finite morphism; restriction to a dense open. Same chapter, the lemma transporting the statements along a finite morphism, and the lemma comparing the statements for a sheaf and for its extension by zero from a dense open. Supplies the transport to X and the final comparison of the two sheaves.
\end{itemize}\medskip
\ifdefined\hypertarget\hypertarget{sourcepointer-1}{}\fi
\noindent\textit{Source reference.} Original reference 06AP: sources/stacks-etale-cohomology/pinned-original-source.tex, source lines 15365--15365. Prerequisite pointer invoked at proof line(s) 15570; it names a standard fact of the same chapter that the counted argument uses, and every step of the counted proof built on it is retained verbatim in the delivered body.\par\medskip
\ifdefined\hypertarget\hypertarget{sourcepointer-2}{}\fi
\noindent\textit{Source reference.} Original reference 0A53: sources/stacks-etale-cohomology/pinned-original-source.tex, source lines 15335--15335. Prerequisite pointer invoked at proof line(s) 15570; it names a standard fact of the same chapter that the counted argument uses, and every step of the counted proof built on it is retained verbatim in the delivered body.\par\medskip
\ifdefined\hypertarget\hypertarget{sourcepointer-3}{}\fi
\noindent\textit{Source reference.} Original reference 03SG: sources/stacks-etale-cohomology/pinned-original-source.tex, source lines 15548--15548. Prerequisite pointer invoked at proof line(s) 15575; it names a standard fact of the same chapter that the counted argument uses, and every step of the counted proof built on it is retained verbatim in the delivered body.\par\medskip
\ifdefined\hypertarget\hypertarget{sourcepointer-4}{}\fi
\noindent\textit{Source reference.} Original reference 0A5D: sources/stacks-etale-cohomology/pinned-original-source.tex, source lines 15471--15471. Prerequisite pointer invoked at proof line(s) 15584; it names a standard fact of the same chapter that the counted argument uses, and every step of the counted proof built on it is retained verbatim in the delivered body.\par\medskip
\ifdefined\hypertarget\hypertarget{sourcepointer-5}{}\fi
\noindent\textit{Source reference.} Original reference 02RX: sources/stacks-etale-cohomology/pinned-original-source.tex, source lines 15491--15491. Prerequisite pointer invoked at proof line(s) 15599; it names a standard fact of the same chapter that the counted argument uses, and every step of the counted proof built on it is retained verbatim in the delivered body.\par\medskip
\ifdefined\hypertarget\hypertarget{sourcepointer-6}{}\fi
\noindent\textit{Source reference.} Original reference 0A52: sources/stacks-etale-cohomology/pinned-original-source.tex, source lines 15324--15324. Prerequisite pointer invoked at proof line(s) 15567; it names a standard fact of the same chapter that the counted argument uses, and every step of the counted proof built on it is retained verbatim in the delivered body.\par\medskip
\ifdefined\hypertarget\hypertarget{sourcepointer-7}{}\fi
\noindent\textit{Source reference.} Original lemma 0BXR: sources/stacks-chow/tags.txt, source lines 14989--14989. Prerequisite pointer invoked at proof line(s) 15578; it names a standard fact of a sibling chapter that the counted argument uses, and every step of the counted proof built on it is retained verbatim in the delivered body.\par\medskip
\ifdefined\hypertarget\hypertarget{sourcepointer-8}{}\fi
\noindent\textit{Source reference.} Original lemma 0B8Y: sources/stacks-chow/tags.txt, source lines 14156--14156. Prerequisite pointer invoked at proof line(s) 15580; it names a standard fact of a sibling chapter that the counted argument uses, and every step of the counted proof built on it is retained verbatim in the delivered body.\par\medskip
\begin{lemma}
\label{lemma-somewhat-easier}
In Situation \ref{situation-what-to-prove} assume $X$ reduced.
Let $j : U \to X$ an open immersion. Let $\ell$ be a prime number
and $\mathcal{F} = j_!\underline{\mathbf{Z}/\ell\mathbf{Z}}$.
Then statements (\ref{item-vanishing}) -- (\ref{item-surjective}) hold
for $\mathcal{F}$.
\end{lemma}

\begin{proof}
The difference with Lemma \ref{lemma-even-easier} is that here we do not
assume $X$ is smooth. Let $\nu : X^\nu \to X$ be the normalization
morphism. Then $\nu$ is finite
(Varieties, Lemma \ref{varieties-lemma-normalization-locally-algebraic})
and $X^\nu$ is smooth
(Varieties, Lemma \ref{varieties-lemma-regular-point-on-curve}).
Let $j^\nu : U^\nu \to X^\nu$ be the inverse image of $U$.
By Lemma \ref{lemma-even-easier} the result holds for
$j^\nu_!\underline{\mathbf{Z}/\ell\mathbf{Z}}$.
By Lemma \ref{lemma-finite-pushforward-statements}
the result holds for $\nu_*j^\nu_!\underline{\mathbf{Z}/\ell\mathbf{Z}}$.
In general it won't be true that
$\nu_*j^\nu_!\underline{\mathbf{Z}/\ell\mathbf{Z}}$ is equal to
$j_!\underline{\mathbf{Z}/\ell\mathbf{Z}}$
but we can work around this as follows.
As $X$ is reduced the morphism $\nu : X^\nu \to X$
is an isomorphism over a dense open $j' : X' \to X$
(Varieties, Lemma \ref{varieties-lemma-normalization-locally-algebraic}).
Over this open we have agreement
$$
(j')^{-1}(\nu_*j^\nu_!\underline{\mathbf{Z}/\ell\mathbf{Z}})
=
(j')^{-1}(j_!\underline{\mathbf{Z}/\ell\mathbf{Z}})
$$
Using Lemma \ref{lemma-restrict-to-open}
twice for $j' : X' \to X$ and the sheaves above
we conclude.
\end{proof}

\end{document}
